% 357_SM_Teilchen_GF27_Zusatz_En_ch.tex
% Johann Pascher, 7 September 2026
% Standard Model quantum numbers from GF(27)* — certificate status (Doc. 357)
\providecommand{\BM}{\textbf{[B]}}
\providecommand{\KM}{\textbf{[K]}}
\providecommand{\SM}{\textbf{[S]}}
\providecommand{\XM}{\textbf{[X]}}

\phantomsection
\addcontentsline{toc}{chapter}{%
  Standard Model Quantum Numbers from GF$(27)^*$ --- Certificate Status (Doc.~357)}

\section*{Status markers}
\begin{center}
\begin{tabular}{@{}lp{10.5cm}@{}}
\textbf{[B]} & \emph{Proved} --- algebraically established. \\[3pt]
\textbf{[K]} & \emph{Confirmed} --- numerically verified. \\[3pt]
\textbf{[S]} & \emph{Speculative} --- open, not derived. \\[3pt]
\textbf{[X]} & \emph{Experimental} --- established; FFGFT derivation open. \\
\end{tabular}
\end{center}

\section*{Description}

This table (Doc.~357) complements the particle diagram (Doc.~356) and
documents the certificate status of all Standard Model quantum numbers
derived algebraically from the multiplicative group GF$(27)^* =
\mathrm{GF}(3^3)^*$.
All closed derivations are without free parameters.

\subsection*{Structure of Doc.~356 (particle diagram)}

Each fermion is shown as a concentric ring diagram. Rings from inside out:
\textbf{white inner circle} = symbol, $Q$, $Y$;
\textbf{colour ring} (quarks only) = $r/g/b$ colour charge \BM\ Doc.~346
--- innermost because gluons confine colour;
\textbf{middle ring} = chirality: left violet $= S_d =$ left-handed,
right orange $= S_u =$ right-handed \BM\ Doc.~349;
\textbf{outer ring} = Galois orbit colour \BM\ Doc.~348;
\textbf{red dashed box} = SU$(2)_L$ doublet pair;
\textbf{black dashed circle} = Generation~3.

\section{Certificate Status}

\subsection*{Fully proved \BM}

\begin{description}

\item[Electric charge $Q$] (Docs.~346, 347 \BM)\\
$Q\in\{0,-1,+2/3,-1/3\}$ from QR-bit $+$ $N\bmod3$.
All four values without free parameters from the Galois classification.

\item[Hypercharge $Y$] (Doc.~347 Theorem~A \BM)\\
$Y = -1 + (4/3)\,N_\text{QR}$; $N_\text{QR}$ = Legendre symbol in GF$(27)^*$.

\item[Weak isospin $I_3$] (Doc.~347 Theorem~B \BM)\\
$I_3 = \pm 1/2$ from the $jE_7$-projector (Su/Sd separation).

\item[Gell-Mann--Nishijima $Q = I_3 + Y/2$] (Doc.~347 Theorem~C \BM)\\
Fully algebraic from the Galois structure without free parameter.

\item[Colour charge $r/g/b$] (Doc.~346 Theorem~C \BM)\\
$Z_3$ charge operator $N\bmod3$ in GF$(3)$; three colours = three elements.

\item[Chirality $S_u\,/\,S_d$] (Doc.~349 \BM)\\
$S_u$ = even Cl$(6)$ grades = right-handed;
$S_d$ = odd = left-handed.

\item[$SU(2)_L$ doublet structure] (Doc.~349 Theorem~D \BM)\\
Doublet = $S_d$-sector (odd grades). Right-handed neutrino $\nu_R$ structurally absent.

\item[3 generations] (Doc.~348 Theorem~A \BM)\\
Three Frobenius cycles of root orbits in GF$(27)^*$; order~3 forced.

\item[Permutation matrix quarks $\times$ leptons] (Doc.~348 Theorem~B \BM)\\
Trace bilinear form GF$(27)$/GF$(3)$: $f_3\times f_4$ structure.

\item[Coupling weights $1/\sqrt{2}$] (Docs.~348 Theorem~C and 353 \BM)\\
$f_1\times f_3$ bilinear form; implies Koide amplitude $d/c=\sqrt{2}$ \BM.

\item[Koide amplitude $d/c=\sqrt{2}$] (Doc.~353 \BM)\\
$Z_3$ equidistribution forces $d/c=\sqrt{2}$, equivalent to $Q=2/3$.
Closed 7~September~2026 (formerly \SM).

\item[8 gluons] (Doc.~339 \BM)\\
$4\times Z_{13}$-orbits $\times$ $2$ $Z_2$-parities; number structurally forced.

\end{description}

\subsection*{Numerically confirmed \KM}

\begin{description}

\item[Electroweak group $Z_2$] (Doc.~336 \KM)\\
$\mathrm{Gal}(\mathrm{GF}(9)/\mathrm{GF}(3)) = Z_2$. Numerically confirmed.

\end{description}

\subsection*{Open / not yet closed \SM}

\begin{description}

\item[Generation assignment (ordering)] (Doc.~346 \SM)\\
Which orbit element corresponds to Generation 1/2/3?
The three orbits are algebraically equivalent; assignment not derived.

\item[CKM/PMNS angles (complete)] (Doc.~348 Theorem~D \SM)\\
Matrix structure established (Theorems B,C \BM);
numerical values require further algebraic work.

\end{description}

\subsection*{Experimentally established, FFGFT derivation open \XM}

\begin{description}

\item[Higgs boson $H^0$] (Doc.~019 \XM)\\
$\tilde{T}\cdot m=1$ replaces Higgs mechanism (Doc.~019).
Boson experimentally established; formal FFGFT derivation open.

\end{description}
